\scanpage{28}
% archive foliation: 27
En vertu de \unsure{tout ce} qui suit,
si $A$ est la sous-catégorie pleine
de $C$ engendrée par les $\Phi(i,n)$
et les $\Psi^i$, alors $C$ est connue
à équiv. près (\unsure{déf. non unique par $A$|déf. can. unique par $A$})
si on connaît $A$ avec sa structure
$k$-additive. Or c'est \inserted{elle} qu'on vient
d'\unsure{être déterminée|être déterminer} (on a omis
de décrire la composition des
morphismes, qui est évidente...).
On trouve donc que $C\simeq$
\unsure{groupes|grps} $\frac12$-simpliciaux ordinaires
str. $\partial$-parfaits ....

La chose restante à faire
est donc déterminer la $\otimes$-
structure et les opérations $\Gamma^i$.
